\documentclass[11pt,a4paper]{article}

\usepackage{graphicx}
\usepackage{xcolor}
\usepackage[left=2.5cm, right=2.5cm, top=2.5cm, bottom=2.5cm]{geometry}
\usepackage{amsmath,amssymb,amsthm}
\usepackage{fontspec}
\setmainfont{Liberation Serif}
\setsansfont{Liberation Sans}
\setmonofont{Liberation Mono}[Scale=0.86]
\usepackage{booktabs}
\usepackage{tabularx}
\newcolumntype{Y}{>{\raggedright\arraybackslash}X}
\usepackage{colortbl}
\usepackage{adjustbox}
\usepackage{enumitem}
\usepackage{tcolorbox}
\definecolor{linkblue}{HTML}{1e3a5f}
\definecolor{accgreen}{HTML}{0e7c66}
\definecolor{warnamber}{HTML}{b45309}
\definecolor{rowlight}{HTML}{f0f5fa}
\usepackage[
    colorlinks=true,
    linkcolor=linkblue,
    citecolor=darkgray,
    urlcolor=linkblue,
    bookmarks=true,
    bookmarksnumbered=true,
    unicode=true
]{hyperref}

\widowpenalty=10000
\clubpenalty=10000
\emergencystretch=1.8em
\tolerance=2500

\newtcolorbox{honestybox}[1][]{colback=warnamber!6, colframe=warnamber!80,
  title={\textbf{Honesty:} #1}, fonttitle=\small, boxrule=0.7pt, arc=2pt}
\newtcolorbox{findbox}[1][]{colback=accgreen!7, colframe=accgreen!80,
  title={\textbf{Finding:} #1}, fonttitle=\small, boxrule=0.7pt, arc=2pt}

\hypersetup{
    pdftitle={Open Problems of the Choptuik--Isaev Framework: Numerical Experiments},
    pdfauthor={Super Z (Z.ai)},
    pdfsubject={Sh.3(iii) bridge, DSI-4, Hurwitz tower, QNM, QCD},
    pdfkeywords={Choptuik, Choptyuk, Klein quartic, Leavitt, DSI, Hurwitz, QNM, GUE}
}

\begin{document}

\tableofcontents
\newpage

% ==========================================================================
\section{Introduction and problem statements}\label{sec:intro}
% ==========================================================================

This report documents the first full computational cycle on the \emph{open
problems} of the Choptuik--Isaev framework developed in the monographs of the
\texttt{choptuik\_ac\_bc} repository. Recall the chain that closes the
Choptyuk problem: the spectral gap on the Klein quartic is described by the
unified formula
\begin{equation}\label{eq:ch}
\Delta_{\mathrm{Ch}} \;=\; \lambda_1(D^2_{\sigma_0}) + \frac{\delta_C^2}{2} - \frac{\delta_C^5}{22},
\qquad \lambda_1(D^2_{\sigma_0}) = 3.338,\quad \delta_C = \frac{\pi}{7},
\end{equation}
which reproduces the observed value $\Delta = 3.443$ within $0.15\%$; the
framework constant $b_{\mathrm{Ch}} = 1-\cos(2\pi/7) \approx 0.3765$ matches
the Choptuik critical collapse exponent within $0.13\%$. Every number below was
produced by an independent run of five experiments
(\texttt{choptuik\_experiments/}), reproducible with a single command
\texttt{python3 run\_all.py}; random seeds are fixed and the results are
deterministic.

The report covers the two directions agreed at the outset. First, the \emph{open
problems}: the Sh.3(iii) bridge --- the only analytic link of the verification
appendix ``Audit and Transfer'' that remains a statement rather than a result
(Experiment~\ref{sec:sh3}); and the DSI-4 observation $c_{K3} = 27/672$, whose
status we strengthen with continued-fraction arithmetic
(Experiment~\ref{sec:dsi4}). Second, \emph{framework extensions}: the Hurwitz
surfaces tower with an audit of the repository data
(Experiment~\ref{sec:hurwitz}), an extended QNM catalogue for LIGO/Virgo with
the overtone ladder (Experiment~\ref{sec:qnm}), and the QCD bridge with an
$N$-scaling test and an extended $\kappa_T$ sweep
(Experiment~\ref{sec:qcd}).

Three new facts emerged, flagged as findings D1, E1 and ``$9/224$ is a
convergent''; each is accompanied by an explicit status (observation / check /
hypothesis) in the spirit of the honesty style adopted in
\texttt{audit\_transfer/}. All tables and figures are generated by code and
refresh automatically when the report is rebuilt.

% ==========================================================================
\section{Experiment 1. The Sh.3(iii) bridge: obstruction ladder}\label{sec:sh3}
% ==========================================================================

\subsection{Statement}

Lemma Sh.3 of the ``Audit and Transfer'' appendix has three parts: (i) the
sharp trace theorem $F \ge 5n/7$ --- proved for all $n$; (ii) the construction
$X_1X_1^* = X_2X_2^* = (4/7)I$ --- the exact equality case; (iii) the
\emph{bridge}: soficity of the group $G = U(L_{F_2}(1,2))$ would induce pairs
$X_1, X_2 \in M_{n_k}(\mathbb{C})$ with per-dimension defect
\begin{equation}\label{eq:defect}
e_k \;=\; \frac{F(X_1,X_2)}{n_k} \;\le\; \eta(\varepsilon_k) \;\to\; 0,
\end{equation}
where
\begin{equation}\label{eq:Fdef}
F(X_1,X_2) = \|X_1^*X_1 - I\|^2 + \|X_2^*X_2 - I\|^2 + \|X_1^*X_2\|^2 + \|X_1X_1^* + X_2X_2^* - I\|^2
\end{equation}
is the Kuncevich defect (Frobenius norm), with universal amplification
$\eta(\varepsilon) = C\varepsilon^{1/2}$. Since (i) gives $e \ge 5/7$ in every
$M_n$, such a bridge would contradict the theorem --- that is the planned proof
of non-soficity of $G$ (the Sh1--Sh4 chain). Part (iii) remains an \emph{exact
statement}. Our experiment asks a practical question: do \emph{natural}
transfers of permutation data into matrices exhibit a decreasing defect? We
build an ``obstruction ladder'' of five families and measure $e$ in each.

\subsection{Methods and results}

\textbf{(A) The theorem at larger $n$.} The identity $F = G(P,Q)$ with
$P = X_1X_1^*$, $Q = X_2X_2^*$ was verified on $300$ random pairs in each
dimension $n = 7,\dots,12$ (maximum relative error at machine level); the Haar
bound $G((\mathrm{tr}P/n)I, (\mathrm{tr}Q/n)I) \le G(P,Q)$ was never violated.
L-BFGS minimization with $15$ restarts yields $|F_{\min} - 5n/7| \le
5.9\cdot10^{-8}$ at $n = 12$ and machine zero for $n \le 10$ --- the sharp
theorem is confirmed beyond the range $n \le 6$ checked in
\texttt{audit\_transfer}.

\textbf{(B--D) Permutation families.} For permutation matrices $P_1, P_2$ the
defect is exact: $F = 2n$, i.e. $e = 2$ (confirmed on $50$ pairs for
$n = 16, 64, 128$ with error at most $5.7\cdot10^{-14}$). The scaled family
$X_i = cP_i$ gives
\begin{equation}\label{eq:scaled}
\frac{F}{n} = 7t^2 - 8t + 3, \qquad t = c^2,
\end{equation}
with minimum exactly $5/7$ at $t = 4/7$ --- the equality case of theorem (i);
the two-parameter family $X_1 = aP_1$, $X_2 = bP_2$ reproduces the same scalar
problem $h(u,v)$ with saddle point $u = v = 4/7$; numerically we obtained
$F/n = 0.7142857143$. This is a closed check: permutations are unitary, so
their scaling is exhausted by the theorem.

\textbf{(E--F) The Cuntz tree.} The truncated shifts $S_1, S_2$ of the binary
tree of depth $k$ ($n = 2^{k+1}-1$) give \emph{exactly}
\begin{equation}\label{eq:tree}
F(S_1, S_2) = n + 2, \qquad e = 1 + \frac{2}{n} \;\to\; 1,
\end{equation}
because the ranges $S_1S_1^*$ and $S_2S_2^*$ partition the depth-$\ge 1$
subspace without overlap ($S_1S_1^* + S_2S_2^* = P(\text{depth} \ge 1)$), so
the last term of (\ref{eq:defect}) equals $\|-P_{\mathrm{root}}\|^2 = 1$ and the
first two equal $2^k$ each (the tree boundary). The weighted family
$X_i = c_iS_i$ gives exactly $F = (n-1)\big[(c_1^2-1)^2 + (c_2^2-1)^2\big] + n + 2$
(verified to $4.5\cdot10^{-13}$); the minimum is attained at $c = (1,1)$ ---
\emph{no weighting improves the truncation}. The tree approximates the Leavitt
relations qualitatively better than permutations ($e \to 1$ versus $e = 2$),
but it neither reaches the floor $5/7$ nor tends to zero.

\textbf{(G) Sofic data.} Simulating condition C1 of Lemma Sh.3, we build pairs
of permutations with fixed-point fraction $\varepsilon$ and apply the natural
transfer $X_i = \sqrt{4/7}\,P_i$. The defect \emph{grows} with $\varepsilon$:
$e(0) = 5/7$, $e(0.01) = 0.745$, $e(0.05) = 0.843$, $e(0.20) = 1.140$. The law
$\eta(\varepsilon) = C\varepsilon^{1/2}$ is not observed in this family ---
on the contrary, the minimum of the defect is attained at $\varepsilon = 0$.

\begin{table}[htbp]
\centering\small
\caption{Obstruction ladder of the Sh.3(iii) bridge: per-dimension defect by
family.}
\label{tab:ladder}
\begin{tabularx}{\textwidth}{lccY}
\toprule
Family & $e = F/n$ & Status & Comment \\
\midrule
Scaled unitaries $c\,U_i$ & $5/7 = 0.7143$ & theorem (equality) & does not approximate the Leavitt relations \\
Cuntz-tree truncation & $1 + 2/n \to 1$ & exact, (\ref{eq:tree}) & best ``honest'' family \\
Permutation models & $2$ & exact & ranges overlap twice \\
Sofic data (transfer) & $\ge 5/7$, grows with $\varepsilon$ & numerical, $n=256$ & no $\sqrt{\varepsilon}$ law \\
\bottomrule
\end{tabularx}
\end{table}

\begin{figure}[htbp]
\centering
\includegraphics[width=0.86\textwidth]{figs_en/fig_sh3_ladder.png}
\caption{Obstruction ladder: per-dimension defect $e = F/n$ for three families
against the sharp floor of the trace theorem.}
\label{fig:ladder}
\end{figure}

\begin{figure}[htbp]
\centering
\includegraphics[width=0.86\textwidth]{figs_en/fig_sh3_hamming.png}
\caption{Experiment 1G: defect of the natural transfer on sofic data with
Hamming error $\varepsilon$; dashed --- the hypothetical law
$\eta(\varepsilon) = C\sqrt{\varepsilon}$.}
\label{fig:hamming}
\end{figure}

\subsection{Conclusion of Experiment 1}

Every constructive family available pays a per-dimension defect of at least
$1$ (tree) or $2$ (permutations); the floor $5/7$ is reachable only by the
trivial scaling of unitaries, which does \emph{not} approximate the Leavitt
relations. Natural transfers exhibit neither decreasing defect nor a
$\sqrt{\varepsilon}$ law. This is a numerical argument that the Sh.3(iii)
bridge is a genuinely non-trivial analytic programme --- and that its ultimate
impossibility would constitute the proof of non-soficity of
$U(L_{F_2}(1,2))$ along the Sh1--Sh4 chain.

\begin{honestybox}[status of the results]
Proved here exactly: $F = n+2$ for the tree truncation (term-by-term proof),
$F = 2n$ for permutations, formula (\ref{eq:scaled}) and the minimum
$(u,v) = (4/7,4/7)$. Verified numerically: the theorem at $n \le 12$ and the
growth of $e(\varepsilon)$. \emph{Not} claimed: non-existence of the bridge
(an analytic statement); the experiment only shows that standard constructions
do not deliver it.
\end{honestybox}

% ==========================================================================
\section{Experiment 2. DSI-4: the observation $c_{K3} = 27/672$}\label{sec:dsi4}
% ==========================================================================

\subsection{Statement}

The amplitude of the log-periodic modulation of the critical cascade,
$c_{K3} = 0.04017757639214903$ (the repository record), is derived from
DSI renormalization at leading order as $b_{\mathrm{Ch}}(22) =
1 - \cos(2\pi/22) = 0.0405070$ (the $+0.82\%$ remainder is explained by the
window systematics of the measurement). Observation DSI-4 records that the
integer combination
\begin{equation}\label{eq:dsi4}
c_{K3} \;=\; \frac{3^3}{2^2\cdot|\mathrm{PSL}(2,7)|} \;=\; \frac{27}{672} \;=\; \frac{9}{224} \;=\; 0.040178571\ldots
\end{equation}
reproduces the measured value within $+0.0025\%$ --- three orders of magnitude
better than the structural alternatives. Status: observation, not a proof. Our
task is to strengthen or weaken this status with independent checks.

\subsection{Results}

\textbf{Continued fraction.} The continued-fraction expansion of the measured
$c_{K3}$ is $[0; 24, 1, 8, 19, 1, \dots]$, whose convergents are $1/24$,
$1/25$, \emph{$9/224$}, $172/4281$, \dots. Hence $9/224 = 27/672$ is a
\emph{convergent of the measured value}: it is not selected from an infinite
pool of rational approximations but arises at the third step of the continued
fraction algorithm. The Gaussian quality measure $|c - p/q|\cdot q^2 = 0.0499
\ll 0.5$ --- an exceptional approximation. Among all $p/q$ with $q \le 224$
exactly one fraction comes as close (namely $9/224$ itself), whereas for a
``random'' target the expected count is $\approx 0.05$.

\textbf{Structural decompositions.} The same number reads in four independent
ways: $3^3/(2^2\cdot168)$ --- the cubic spinor structure $B^3 = -I$ over the
quadratic contribution $A^2 = -I$, normalized by the automorphism group order;
$9/224 = 3^2/(2^5\cdot 7)$; $(N-1)/(24N)$ with $N = 28$ --- the Hilbert-space
dimension of the framework over the $24$ heptagons of the $\{7,3\}$ tiling of
the Klein quartic; and $27/(12\cdot 56)$ over the $56$ triangles of the same
tiling. Every decomposition yields exactly $0.040178571\ldots$

\textbf{Candidate table.} Among the thirteen candidates checked, only $27/672$
is better than $0.4\%$; the unstructured $1/25$ lags two orders behind.
Window-compatibility: the ratio measured/candidate $= 0.99998$ for $27/672$
requires a nearly zero window deficit, whereas $b_{\mathrm{Ch}}(22)$ requires
$0.81\%$ (inside the observed band $[0.987, 0.999]$). Both readings are
compatible with the measurement protocol; $27/672$ is, moreover, $300$ times
more accurate.

\begin{table}[htbp]
\centering\small
\caption{Best candidates for the amplitude $c_{K3}$ (full table in JSON).}
\label{tab:cands}
\begin{tabularx}{\textwidth}{lccY}
\toprule
Candidate & Value & Dev., \% & Status \\
\midrule
$27/672 = 3^3/(2^2\cdot|\mathrm{PSL}(2,7)|)$ & $0.040178571$ & $+0.0025$ & DSI-4 observation, strengthened \\
$1/25$ & $0.040000000$ & $-0.442$ & unstructured \\
$b_{\mathrm{Ch}}(22(1+\gamma))$ (braked RG) & $0.040359588$ & $+0.452$ & DSI-2 \\
$b_{\mathrm{Ch}}(22)$ & $0.040507026$ & $+0.820$ & DSI-1 \\
$1-\cos(2\pi/24)$ & $0.034074174$ & $-15.19$ & rejected \\
\bottomrule
\end{tabularx}
\end{table}

\begin{findbox}[$9/224$ is a convergent]
The DSI-4 observation is strengthened: $27/672$ is not merely a ``good
approximation'' but a member of the convergent sequence of the measured value
with Gaussian measure $0.0499$. The analytic derivation of $3^3/(2^2\cdot168)$
from DSI renormalization remains open; the problem now has an exact numerical
target.
\end{findbox}

\begin{figure}[htbp]
\centering
\includegraphics[width=0.92\textwidth]{figs_en/fig_dsi4_candidates.png}
\caption{Structural candidates for the amplitude $c_{K3}$: absolute deviation
from the measured value (log scale).}
\label{fig:cands}
\end{figure}

% ==========================================================================
\section{Experiment 3. The Hurwitz surfaces tower}\label{sec:hurwitz}
% ==========================================================================

\subsection{Audit of the repository data: finding D1}

The monograph (Part 7 of \texttt{verify\_enhanced}) uses the tower
$\{$Klein $g{=}3$, $n{=}7$; Macbeath $g{=}7$, $n{=}9$; Hurwitz\_3 $g{=}14$,
$n{=}11\}$. We checked every row against two mandatory relations for a Hurwitz
surface: $|\mathrm{Aut}| = 84(g-1)$ and $n \mid |\mathrm{Aut}|$ (an element of
order $n$). The first two rows are consistent: $\mathrm{PSL}(2,7)$,
$168 = 84\cdot 2$, $7\mid 168$; $\mathrm{PSL}(2,8)$, $504 = 84\cdot 6$,
$9\mid 504$. The third row is \textbf{inconsistent}: $84(14-1) = 1092 =
2^2\cdot 3\cdot 7\cdot 13$ is not divisible by $11$, so a group of order
$1092$ has no elements of order $11$, and the signature $(2,3,11)$ is
impossible for genus $14$.

\begin{findbox}[D1: repository tower inconsistency]
For $g = 14$ the actual Hurwitz surface is the quotient by $\mathrm{PSL}(2,13)$,
$|\mathrm{PSL}(2,13)| = 1092$, with signature $(2,3,13)$. In the corrected
tower $\delta = \pi/13$, not $\pi/11$; all framework computations for this
point must be recomputed. We consider this a data editing error (not an error
in the mathematics of the corrections).
\end{findbox}

\subsection{The corrected tower and computations}

We build the corrected tower of the four smallest Hurwitz surfaces, including
the \emph{first triplet} of genus $17$ ($|G| = 1344 = 84\cdot 16$, signature
$(2,3,8)$) --- the nearest point after Klein where the tower branches. All
checks $84(g-1) = |\mathrm{Aut}|$, $n \mid |\mathrm{Aut}|$ and hyperbolicity
$1/2 + 1/3 + 1/n < 1$ pass. For every surface we compute the framework phase
quantities $\delta = \pi/n$, $\gamma = \delta^4/b_2$,
$\delta_{\mathrm{eff}} = \delta^5/b_2$ under the hypothetical pattern
$b_2(k) = 2b_2(k-1) + 2 = 3\cdot 2^{k+3} - 2$, which reproduces the three
repository values $(22, 46, 94)$ and extends them to the fourth point ($190$).

\begin{table}[htbp]
\centering\small
\caption{The corrected Hurwitz tower and framework quantities. $b_2(k)$ is the
pattern hypothesis (topological only for Klein: $b_2(K3) = 22$).}
\label{tab:tower}
\begin{tabular}{llcclll}
\toprule
Surface & Group & $g$ & $n$ & $|\mathrm{Aut}|$ & $\gamma$ & $\delta_{\mathrm{eff}}$ \\
\midrule
Klein & PSL(2,7) & 3 & 7 & 168 & $1.844\cdot10^{-3}$ & $8.276\cdot10^{-4}$ \\
Macbeath & PSL(2,8) & 7 & 9 & 504 & $3.228\cdot10^{-4}$ & $1.127\cdot10^{-4}$ \\
PSL(2,13) & PSL(2,13) & 14 & 13 & 1092 & $3.628\cdot10^{-5}$ & $8.768\cdot10^{-6}$ \\
First triplet & $(2,3,8)$, $|G|=1344$ & 17 & 8 & 1344 & $1.252\cdot10^{-4}$ & $4.915\cdot10^{-5}$ \\
\bottomrule
\end{tabular}
\end{table}

The pressure on $\delta_{\mathrm{eff}}$ along the PSL branch is $\times 7.3$
(Klein $\to$ Macbeath) and $\times 12.9$ (Macbeath $\to$ PSL(2,13)); the
branching at the triplet returns part of the amplitude ($n$ drops from $13$ to
$8$). The universal curve $\Delta_{\mathrm{Ch}}(\delta) = 3.338 + \delta^2/2 -
\delta^5/22$ increases monotonically for $\delta \in (0, 1.64)$ (derivative
$\delta(1 - 5\delta^3/22) > 0$), so the ``minimum near $\pi/7$'' in the
monograph should correctly be read as the crossing point of the curve with the
observed value: $\delta_* = 0.4603$ versus $\pi/7 = 0.4488$ (difference
$0.026$). The points $\pi/n$ of the corrected tower give
$\Delta_{\mathrm{Ch}}(\pi/13) = 3.3672$ --- a framework prediction for the
genus-14 spectral quantity, to be checked once $\lambda_1$ is known for that
surface.

\begin{honestybox}[the hypothetical $b_2(k)$]
For non-Klein surfaces of the tower the value $b_2$ is not derived from
topology: only $b_2 = 22 = b_2(K3)$ for Klein is topological. The pattern
$b_2(k) = 3\cdot 2^{k+3}-2$ is an extrapolation of three repository points;
all numbers in Table~\ref{tab:tower} should be read as a \emph{structural
hypothesis} of the framework. The formal continuation beyond the fourth point
($n = 10, 12, 14$ with $b_2 = 382, 766, 1534$) is not tied to real surfaces
and is listed only as a pattern continuation.
\end{honestybox}

\begin{figure}[htbp]
\centering
\includegraphics[width=0.88\textwidth]{figs_en/fig_hurwitz_tower.png}
\caption{The corrected Hurwitz tower: braking $\gamma$ and effective phase
$\delta_{\mathrm{eff}}$ across surfaces (log scale).}
\label{fig:tower}
\end{figure}

% ==========================================================================
\section{Experiment 4. QNM: the LIGO/Virgo catalogue and overtones}\label{sec:qnm}
% ==========================================================================

\subsection{Statement and method}

Framework v2.0 predicts a multiplicative correction to black-hole
quasi-normal mode frequencies:
\begin{equation}\label{eq:qnm}
f_{\mathrm{corr}} = f\left(1 - \frac{\delta_{\mathrm{eff}}}{\pi^2}\right)
= 0.999916143\,f, \qquad
\frac{\Delta f}{f} = -8.3857\cdot 10^{-5}.
\end{equation}
We extend the repository catalogue from four to ten events. For the four
``anchor'' events we use the ringdown frequencies measured in the repository;
for six new events the $l = m = 2$ fundamental frequency is computed from the
Berti--Cardoso--Will fit with final masses and spins from the LVK papers:
$M\omega_R = 1.5251 - 1.1568(1-a)^{0.1292}$. For the anchors both estimates
are shown side by side --- the $8$--$13\%$ spread reflects the accuracy of the
fit and of the ringdown extraction method, not a data error.

\begin{table}[htbp]
\centering\small
\caption{The ten-event catalogue: predicted shifts $\Delta f$. For anchors
($\ast$) $f$ is measured (repository), otherwise --- Berti fit.}
\label{tab:qnm}
\begin{tabular}{lrrrrr}
\toprule
Event & $M_f$, $M_\odot$ & $a_f$ & $f$, Hz & $\Delta f$, Hz & $f_{\mathrm{corr}}$, Hz \\
\midrule
GW150914$^*$ & 62.3 & 0.67 & 251.0 & $-0.0210$ & 250.979 \\
GW151226 & 20.5 & 0.74 & 871.8 & $-0.0731$ & 871.710 \\
GW170104$^*$ & 48.7 & 0.65 & 314.0 & $-0.0263$ & 313.974 \\
GW170608 & 17.8 & 0.69 & 963.5 & $-0.0808$ & 963.419 \\
GW170729 & 59.1 & 0.81 & 323.5 & $-0.0271$ & 323.473 \\
GW170809 & 35.4 & 0.70 & 488.3 & $-0.0409$ & 488.259 \\
GW170814$^*$ & 53.4 & 0.70 & 286.0 & $-0.0240$ & 285.976 \\
GW170818 & 35.6 & 0.67 & 474.4 & $-0.0398$ & 474.360 \\
GW190412 & 40.3 & 0.87 & 510.2 & $-0.0428$ & 510.157 \\
GW190521$^*$ & 142.0 & 0.72 & 110.0 & $-0.0092$ & 109.991 \\
\bottomrule
\end{tabular}
\end{table}

\subsection{Detectability and overtones}

Because the correction (\ref{eq:qnm}) is multiplicative, the ratio
$|\Delta f|/\sigma$ is determined by the relative detector precision alone:
$0.004\sigma$ at the current $2\%$, $0.084\sigma$ for the Einstein Telescope
($0.1\%$) and $0.84\sigma$ for Cosmic Explorer ($0.01\%$). The $1\sigma$
threshold requires precision $0.0084\%$; combining $N$ independent events
yields a $\sqrt{N}$ gain, so a catalogue of ten events already crosses the
threshold in aggregate at CE-class detectors. This refines the repository
claim ``near-detectable'' and makes the prediction falsifiable on the horizon
of third-generation detectors.

The overtone ladder for GW150914 is built in the Schwarzschild limit from the
reference values $M\omega_n$ (Berti et al., 2006): $f_0..f_4 = 251,\ 394,\
537,\ 681,\ 825$ Hz with damping times $\tau_0..\tau_4 = 3.45,\ 1.57,\ 1.02,\
0.76,\ 0.60$ ms. The multiplicativity of the correction yields a falsifiable
prediction: $\Delta f_n / f_n = \mathrm{const} = 8.3857\cdot10^{-5}$ for
\emph{all} overtones --- the relative spacings $f_n/f_0$ are unchanged.
Overtone spectroscopy detecting disproportionate shifts would falsify the
framework; proportional shifts with the stated constant would confirm it.

\begin{honestybox}[status of the QNM numbers]
For non-anchor events $f$ is a GR prediction from $(M_f, a_f)$, not a
measurement; the correction is applied to the GR value. The overtones are in
the Schwarzschild limit: at $a = 0.67$ the frequencies shift by $O(1-a) \sim
30\%$; the ladder is used as a scaling scheme, not as an exact Kerr table.
\end{honestybox}

\begin{figure}[htbp]
\centering
\includegraphics[width=0.88\textwidth]{figs_en/fig_qnm_overtones.png}
\caption{The GW150914 overtone ladder ($l = m = 2$, Schwarzschild limit,
scaled to the measured $f_0 = 251$ Hz): frequencies $f_n$ and shifts
$\Delta f_n$.}
\label{fig:overtones}
\end{figure}

% ==========================================================================
\section{Experiment 5. The QCD bridge: $N$-scaling and $\kappa_T$ sweep}\label{sec:qcd}
% ==========================================================================

\subsection{Finding E1: the trace of the operator $O_\chi$}

The work formula of the framework, $\bar\theta_{\mathrm{eff}} = \delta_C\,
N\langle\lambda\rangle \mathcal{S}_{\mathrm{GUE}}$, relies on the spectral
symmetry of the GUE class: $\langle\lambda\rangle = 0$ in the continuum. We
built the \emph{tower} $O_\chi(N = 28k)$ of $k$ blocks $(Q_{K3} \oplus M_F)$
with the full $T$-breaking matrix $V_T$ and found: the trivial trace of the
core is $\mathrm{tr}(Q_{K3}\oplus M_F) = 32$ (the $E_8$ Cartan diagonals equal
$2$), so the raw mean $\langle\lambda\rangle = 32/28 = 1.143$ is a
\emph{constant} along the tower, and the symmetry argument does not apply to
the raw operator. The fix is one line and physically mandatory: $\bar\theta$
couples to the \emph{fluctuation} of the topological charge, not to its
constant offset, so the operator must be centered: $O_\chi \to O_\chi -
(\mathrm{tr}O_\chi/N)I$. After centering, the fluctuation of the sample mean
decays as a power law with measured slope $-0.99$, i.e.
$|\langle\lambda\rangle| \sim 1/N$ --- even faster than the announced
$1/\sqrt{N}$.

\begin{findbox}[E1: centering of $O_\chi$]
In the repository construction of $O_\chi$ the trace $32$ is not removed. The
work formula requires centering; after it, $\langle\lambda\rangle \to 0$ with
slope $-0.99$ in $\log N$ (tower $N = 28..448$, $M = 48$ realizations per
point). This is an editorial fix of the construction that strengthens the
framework's conclusion.
\end{findbox}

\subsection{The $\kappa_T$ sweep: the GUE crossover}

An extended scan $\kappa_T \in [0, 16]$ at $N = 28$ ($200$ realizations of
$V_T$ per point; GUE/GOE/Poisson reference ensembles of $3000$ samples each)
reproduces and refines the repository picture. At $\kappa_T = 0$ the spectrum
is degenerate (an integer block operator); as the coupling grows, the spacing
statistics cross monotonically from Poisson to GUE: $\bar r$ rises from
$0.045$ to $\approx 0.36$ (GUE reference $0.387$, Poisson $1/3$). The Bayes
factor $\mathrm{BF}(\mathrm{GUE}/\mathrm{Poi})$ reaches the ``decisive'' level
($\ge 100$) at $\kappa_T \approx 4$ and peaks at $\kappa_T = 8.45$
($\mathrm{BF} \approx 1165$), consistent with the lattice estimate
$\kappa_T^{(\mathrm{QCD})} > 2.62$ (95\% CL) under which the repository
interpolates $\mathrm{BF} \ge 99$. In our independent grid we obtain
$\mathrm{BF} = 50$ (strong) at $\kappa_T = 2.62$; the qualitative conclusion
(``at the physical $\kappa_T$ the operator is in the GUE class'') is
reproduced.

\begin{table}[htbp]
\centering\small
\caption{The $\kappa_T$ sweep at $N = 28$ after centering (finding E1): mean
folded ratio and Bayes factor.}
\label{tab:kappa}
\begin{tabular}{rcccl}
\toprule
$\kappa_T$ & $\bar r$ & $\sigma_r$ & BF(GUE/Poi) & Verdict \\
\midrule
0.00 & 0.045 & 0.000 & $0.0$ & degenerate \\
0.50 & 0.243 & 0.020 & $0.0$ & Poisson class \\
1.00 & 0.308 & 0.022 & $0.8$ & crossover \\
1.50 & 0.333 & 0.023 & $15.5$ & STRONG GUE \\
2.62 & 0.344 & 0.022 & $50.3$ & STRONG GUE \\
4.00 & 0.353 & 0.023 & $193.3$ & DECISIVE GUE \\
8.45 & 0.362 & 0.024 & $1164.9$ & DECISIVE GUE \\
16.00 & 0.358 & 0.024 & $501.8$ & DECISIVE GUE \\
\bottomrule
\end{tabular}
\end{table}

\begin{honestybox}[the $V_T$ model and the BF grid]
$V_T$ remains a model matrix of $T$-violation (a Gaussian Hermitian block
scaled by $b_C$) --- a constructive hypothesis of the framework. The BF
numbers depend on the grid and the KDE kernel; comparison with the repository
interpolations ($99$ at $2.62$, $510$ at $8.45$) is qualitative, not numeric.
\end{honestybox}

\begin{figure}[htbp]
\centering
\includegraphics[width=0.86\textwidth]{figs_en/fig_qcd_nscaling.png}
\caption{$N$-scaling of the centered fluctuation of $\langle\lambda\rangle$ on
the $O_\chi$ tower: measured slope $-0.99$ against the $1/\sqrt{N}$ and $1/N$
references.}
\label{fig:nscaling}
\end{figure}

\begin{figure}[htbp]
\centering
\includegraphics[width=0.86\textwidth]{figs_en/fig_qcd_kappa_sweep.png}
\caption{The $\kappa_T$ sweep: Bayes factor BF(GUE/Poisson); the vertical
line is the lattice threshold $\kappa_T > 2.62$ (95\% CL).}
\label{fig:kappa}
\end{figure}

% ==========================================================================
\section{Honesty and limitations}\label{sec:honesty}
% ==========================================================================

Proved analytically (within the experiment): the exact formula $F = n + 2$ for
the Cuntz-tree truncation with a term-by-term proof; the exact formula
$F = (n-1)[(c_1^2-1)^2 + (c_2^2-1)^2] + n + 2$ for the weighted tree and its
minimum at $c = (1,1)$; the formula $F/n = 7t^2 - 8t + 3$ for scaled
permutations with minimum $5/7$ at $t = 4/7$; the chain of equalities
(\ref{eq:dsi4}). Verified numerically with stated accuracy: the sharp trace
theorem at $n \le 12$ (L-BFGS, gap $\le 5.9\cdot10^{-8}$); the
exceptionality of $9/224$ (Gaussian measure $0.0499$; one fraction at
$q \le 224$ against an expectation of $0.05$); the growth of $e(\varepsilon)$
on sofic data; the power law of the centered fluctuation of
$\langle\lambda\rangle$ (slope $-0.99$).

Not claimed, and still open: the non-existence of the Sh.3(iii) bridge (our
result is only the failure of standard constructions to deliver it); the
analytic derivation of $27/672$ from DSI renormalization (the problem now has
an exact target --- the convergent); the derivation of the tower's $b_2(k)$
from topology; the $V_T$ model in the QCD experiments. Findings D1 (the
$n = 11$ inconsistency of the repository tower) and E1 (the trace $32$ in
$O_\chi$) are formulated as editorial fixes of data and construction, not as
refutations of the mathematics of the corrections.

% ==========================================================================
\section{Conclusions and next steps}\label{sec:concl}
% ==========================================================================

The first full cycle of work on the open problems of the framework is
complete. For the Sh.3(iii) bridge we built the obstruction ladder: between
the best ``honest'' family (the Cuntz tree, $e \to 1$) and the sharp floor
($5/7$) a gap remains that is closed neither by permutations, nor by
weighting, nor by sofic data with the natural transfer --- so the analytic
programme (iii) now has a numerical contour of the target construction. For
DSI-4 the observation $27/672$ is strengthened by the convergent status with
Gaussian measure $0.0499$ and four independent structural decompositions. For
the Hurwitz tower we found and fixed a repository data inconsistency (D1) and
built the corrected four-point tower with the prediction
$\Delta_{\mathrm{Ch}}(\pi/13) = 3.3672$. The QNM catalogue is extended to ten
events with an explicit falsification horizon (precision $0.0084\%$,
CE-class detectors) and a falsifiable overtone prediction
$\Delta f_n/f_n = \mathrm{const}$. In the QCD experiment we identified the
mandatory centering fix of $O_\chi$ (E1), after which the $N$-scaling is
confirmed with slope $-0.99$ and the $\kappa_T$ sweep reproduces the GUE
crossover with $\mathrm{BF} \ge 50$ already at the lattice threshold.

The next steps follow from the findings. First, apply the fixes D1 and E1 to
the repository (the Hurwitz\_3 row and the $O_\chi$ construction) and rerun
its verification. Second, attack the analytic target of DSI-4: explain why
DSI renormalization with $\lambda = 22$ singles out the convergent $9/224$;
a natural path is to check whether the measurement window fixes the
denominator $224 = 4\cdot 56$ through the combinatorics of the $\{7,3\}$
tiling. Third, close the gap $[5/7,\,1]$ of the Sh.3(iii) bridge: prove the
lower bound $e \ge 1$ for the class of ``isometric-on-support'' constructions
to which the tree belongs --- this would separate tree truncation from scaled
unitaries and turn the obstruction ladder into a theorem. Finally, the QNM
direction benefits from recomputing the overtones in the Kerr limit
($a = 0.67$) from exact tables, replacing the scaling scheme with a
full-accuracy prediction.

% ==========================================================================
\appendix
\section{Reproducibility}\label{sec:repro}
% ==========================================================================

All experiments were run by the orchestrator \texttt{run\_all.py}
(Python 3.12 with NumPy/SciPy/Matplotlib), which launches the five scripts
with a single command; a full run takes about five minutes. Random seeds are
fixed: 42 --- the QCD reference ensembles;
5000 + 97r + n --- the L-BFGS restarts; 9000 + 131r + k --- the tree;
11000 + \ldots{} --- the sofic data. Two consecutive runs produce
bit-identical JSON files. Results live in the results directory (files
exp*.json), figures in the fig\_ru and fig\_en directories, one set per
report language. Integrity checks passed on the same environment before the
experiments: the repository test suite (30 tests), the audit\_transfer
verification and the full verify\_enhanced verification.

\end{document}
